% ============================================================
% APPENDIX for Submission 281
% "One Flood Substrate, Six Decision Geometries:
%  A Diagnostic Benchmark for Spatial AI [Experiment]"
%
% Purpose: paper-supporting appendix. Substantiates the abstract's two
% load-bearing claims —
%   (1) "the benchmark organizes evaluation around six spatial decision
%        geometries"  -> the protocol (evidence mapping + verdict logic)
%   (2) "per-feature provenance and field-status annotations" + "where the
%        evidence is missing or provisional" -> the substrate (sources,
%        field-status semantics, availability matrix)
%
% Scope discipline: this appendix carries ONLY what the abstract names.
% No certificate gate, action vocabulary, gearbox, oobleck, MoE/forecast
% stack, or controlplane architecture appears here — none is in the
% abstract, so none belongs in an appendix to it.
%
% Integration:
%   - assumes \appendix has been declared in the body
%   - \PH{} is the red placeholder macro (defined defensively below if absent)
%   - requires \usepackage{booktabs} and \usepackage{pifont}
% ============================================================


\section{Decision-Geometry Evaluation Protocol and Substrate}
\label{app:protocol-substrate}

\paragraph{Two distinct ``sixes.''}
The benchmark is organized around six \emph{decision} geometries---the
operational task types a crisis manager performs: prediction, ranking,
clustering, transfer, relational propagation, and allocation. These are
\emph{not} the six measurement proxies of the underlying compatibility
machinery (smooth, spatial, composite, hierarchical, logical, periodic),
which are diagnostic instruments. The two are orthogonal: a kappa proxy
may inform several decision geometries, and no decision geometry maps
one-to-one onto a single proxy. Indeed, as
Table~\ref{app:tab:evidence-map} shows, five of the six decision
geometries are evaluated by \emph{direct} comparison of model outputs
(rank correlations, leave-event-out deltas, ablation uplift); only
\emph{clustering} is evaluated through a registered kappa proxy
(\textit{kappa\_spatial}, via Moran's $I$ on residuals). The evaluation
axis is the decision, not the proxy.

\subsection{Protocol: per-geometry evidence}
\label{app:protocol}
Table~\ref{app:tab:evidence-map} states, for each decision geometry, the
test that supplies its evidence, the diagnostic it relies on, and its
implementation status as of this writing. We report status honestly:
results are not locked, and one geometry's evidence path is specified but
not yet implemented in the training code.

\begin{table*}[t]
\centering
\caption{Decision-geometry evidence map. ``Test'' is the artifact that
  supplies evidence; ``Diagnostic'' is what the test relies on (only
  clustering uses a registered kappa proxy); ``Status'' is the
  implementation state, reported honestly pending data lock. Every
  reported result downstream of this table is a placeholder until runs
  complete.}
\label{app:tab:evidence-map}
\small
\begin{tabular}{@{}p{2.5cm}p{4.3cm}p{3.4cm}p{4.2cm}@{}}
\toprule
\textbf{Decision geometry} & \textbf{Test (evidence artifact)} & \textbf{Diagnostic} & \textbf{Status} \\
\midrule
Prediction &
  R0$\to$R1$\to$R2 representation ladder; fold-level Wilcoxon on paired deltas &
  ladder skill ($R^2$, RMSE) &
  Specified; R0 executing \\[3pt]
Ranking &
  Within-event rank shift R0 vs R1 (Kendall's $\tau$-b, top-$k$ churn, fidelity-to-truth $\Delta$) &
  none --- direct rank comparison &
  Specified (ranking-geometry design); not executed \\[3pt]
Clustering &
  Moran's $I$ on spatial-blocked residuals, per level &
  \textit{kappa\_spatial} (registered) &
  Wired (\texttt{compute\_diagnostics.py} imports \texttt{yrsn.core.kappa.spatial}); not executed \\[3pt]
Transfer &
  Leave-event-out vs.\ spatial-blocked metric ratio &
  none --- direct metric comparison &
  Specified in all training scripts; not executed \\[3pt]
Relational propagation &
  W-matrix ablation (R1-full vs.\ R1-no-wlag vs.\ no-target-lag vs.\ wlag-only) &
  none --- direct uplift comparison (\textit{kappa\_logical} is the conceptual instrument, not computed) &
  \textbf{Specified in DOE; not implemented.} Training script produces R1-full only; ablation-variant predictions pending \\[3pt]
Allocation &
  R2 temporal lead-time features (MRMS/HRRR rainfall timing) &
  none --- direct (\textit{kappa\_periodic} is the conceptual instrument, not computed) &
  \textbf{Partial.} Coverage gap by data vintage (\S\ref{app:substrate}, Table~\ref{app:tab:availability}) \\
\bottomrule
\end{tabular}
\end{table*}

\subsection{Per-geometry sufficiency verdicts}
\label{app:verdicts}
The benchmark's output is not a leaderboard but a \emph{sufficiency
verdict} per geometry: \textsc{supported}, \textsc{provisional}, or
\textsc{insufficient}. The verdict criteria are fixed before results.

\begin{description}
\item[Prediction.] \textsc{Supported} if the ladder beats the mean
  predictor and a random-feature null under spatial-blocked CV
  ($R^2>0$ on $\geq 2$ targets) and R1 improves on R0 at the fold level
  (Wilcoxon $p<0.05$ \emph{and} Cohen's $d>0.2$); else \textsc{provisional}.
\item[Ranking.] With reordering threshold pre-committed at Kendall's
  $\tau\text{-b}<0.90$: \textsc{supported} if context meaningfully reorders
  \emph{and} moves the order toward observed impact
  ($\Delta\text{-fidelity}>0$); \textsc{provisional} if it reorders without
  improving fidelity; context not load-bearing if $\tau\text{-b}\geq 0.90$.
  Because NFIP claims are zero-inflated, the headline is the non-zero
  (claimant) subset when tied pairs exceed half.
\item[Clustering.] \textsc{Supported} if residual Moran's $I$ is
  indistinguishable from zero after the spatial ladder level;
  \textsc{insufficient} if errors remain spatially clustered (a standing
  evidence gap for any spatially-resolved decision).
\item[Transfer.] \textsc{Supported} if leave-event-out skill retains a
  pre-committed fraction of spatial-blocked skill; \textsc{provisional}
  otherwise. Cross-event is reported; cross-\emph{city} is weaker evidence
  and not pooled.
\item[Relational propagation.] \textsc{Supported} if R0$\to$R1 uplift
  survives the no-target-lag ablation; if uplift is entirely the target
  spatial lag, the honest verdict is \textsc{provisional}---``neighbor
  activity predicts,'' not ``relational structure is supported.''
  \emph{This verdict cannot be issued until the ablation variants are
  implemented} (Table~\ref{app:tab:evidence-map}).
\item[Allocation.] \textsc{Supported} for events with temporal coverage;
  \textsc{insufficient (evidence absent by vintage)} for events predating
  the rainfall-timing sources (\S\ref{app:substrate}). The verdict is
  therefore necessarily \emph{partial} across the event panel---a
  demonstration of the benchmark surfacing where a decision lacks support,
  not a defect.
\end{description}

\subsection{Substrate, provenance, and field status}
\label{app:substrate}
The benchmark runs on one multi-event flood substrate at the
(\textit{zcta\_id}, \textit{event}) grain, spanning five metropolitan
regions. The nine modelable cells---the basis of the $n{=}9$ figure used
throughout---are the (scenario, target) pairs with sufficient coverage to
train, enumerated in Table~\ref{app:tab:cells}.

\begin{table}[t]
\centering
\caption{The nine modelable (scenario, target) cells. A cell is modelable
  when the target is populated at the (\textit{zcta\_id}, \textit{event})
  grain with enough non-missing rows to fit under spatial-blocked CV;
  $n$ is rows (ZCTAs $\times$ events). This table is what grounds the
  $n{=}9$ claim.}
\label{app:tab:cells}
\small
\begin{tabular}{@{}llll@{}}
\toprule
\textbf{Scenario} & \textbf{Target} & \textbf{Task} & \textbf{$n$ rows} \\
\midrule
Houston      & \texttt{obs\_nfip\_event\_claims} & regression     & 396 \\
Houston      & \texttt{obs\_has\_311}            & classification & 396 \\
Houston      & \texttt{obs\_has\_hwm}            & classification & 396 \\
SW Florida   & \texttt{obs\_nfip\_event\_claims} & regression     & 606 \\
NYC          & \texttt{obs\_nfip\_event\_claims} & regression     & 422 \\
NYC          & \texttt{obs\_has\_311}            & classification & 422 \\
Riverside    & \texttt{obs\_nfip\_event\_claims} & regression     & 172 \\
New Orleans  & \texttt{obs\_nfip\_event\_claims} & regression     & 264 \\
New Orleans  & \texttt{obs\_has\_hwm}            & classification & 264 \\
\bottomrule
\end{tabular}
\end{table}

Each feature carries provenance and a \emph{field-status}
annotation drawn from a fixed vocabulary, so that an absent value is never
silently treated as a measured zero:

\begin{description}
\item[\texttt{available}.] The source covers this event; a value is present.
\item[\texttt{not\_applicable\_temporal}.] The event predates the source's
  operational start. This is a vintage boundary, \emph{not} a missing
  measurement.
\item[\texttt{not\_applicable\_type}.] The source does not apply to this
  event type (e.g.\ Atlantic best-track for an atmospheric-river event).
\item[\texttt{source\_empty}.] The source covers the event but returned
  zero records.
\end{description}

This annotation is what makes the allocation verdict legible. Allocation
depends on rainfall-timing features (NOAA MRMS Stage~IV; HRRR) and soil
moisture (SMAP). MRMS~v12 is operational from October~2012, HRRR from
September~2014, and SMAP from March~2015. Every event in the panel falls
within coverage \emph{except} the two earliest New Orleans events, which
are marked \texttt{not\_applicable\_temporal} rather than missing
(Table~\ref{app:tab:availability}). For those events, temporal lead-time
sufficiency cannot be assessed, and allocation is reported as
\textsc{insufficient (evidence absent by vintage)}.

\begin{table}[t]
\centering
\caption{Temporal-source coverage for the allocation geometry, shown on
  the New Orleans event series (the pre/post-levee substrate, where the
  vintage boundary bites). \ding{51}\,\texttt{available};
  \ding{55}\,\texttt{not\_applicable\_temporal} (predates source). All
  non-NOLA panel events postdate October~2012 and carry MRMS.}
\label{app:tab:availability}
\small
\begin{tabular}{@{}llccc@{}}
\toprule
\textbf{Event} & \textbf{Date} & \textbf{MRMS} & \textbf{HRRR} & \textbf{SMAP} \\
\midrule
Katrina & 2005-08-29 & \ding{55} & \ding{55} & \ding{55} \\
Isaac   & 2012-08-28 & \ding{55} & \ding{55} & \ding{55} \\
Barry   & 2019-07-11 & \ding{51} & \ding{51} & \ding{51} \\
Ida     & 2021-08-26 & \ding{51} & \ding{51} & \ding{51} \\
\midrule
\multicolumn{2}{@{}l}{Source operational from} & 2012-10 & 2014-09 & 2015-03 \\
\bottomrule
\end{tabular}
\end{table}

\paragraph{Reproducibility note.}
The field-status vocabulary above is the single source of truth consumed
by the data-readiness validator; an event/source pair marked
\texttt{not\_applicable\_temporal} is excluded from the affected geometry's
evidence rather than imputed, so the ``missing or provisional'' verdicts
are determined by the substrate annotation, not by post-hoc judgment.

\subsection{Relational ablation and leakage controls}
\label{app:relational-ablation}
The relational-propagation geometry (\S\ref{sec:geom-relational}) is
evaluated by ablating the R1 W-matrix block. The W-matrix adds eight
spatial-structure features to the R0$\to$R1 step: node degree, mean
neighbor distance, four static neighbor lags, a rainfall lag
(\textit{wlag\_rainfall\_mm}), and the lag of the target itself
(\textit{wlag\_nfip\_claims}). Four feature sets isolate what the spatial
structure contributes:

\begin{table}[h]
\centering
\caption{W-matrix ablation battery. Feature counts are for the Houston
  cell; the design is identical across scenarios.}
\label{app:tab:ablation}
\small
\begin{tabular}{@{}llc@{}}
\toprule
\textbf{Variant} & \textbf{Features} & \textbf{$n_\text{feat}$} \\
\midrule
full          & R0 + hydrology + W-matrix              & 61 \\
no-wlag       & R0 + hydrology only                    & 53 \\
no-target-lag & full minus \textit{wlag\_nfip\_claims} & 60 \\
wlag-only     & R0 + W-matrix only                     & 41 \\
\bottomrule
\end{tabular}
\end{table}

\paragraph{Verdict logic.} The relational verdict rests on
\textbf{no-target-lag}: if the R0$\to$R1 uplift survives removing
\textit{wlag\_nfip\_claims}, spatial structure is supported as evidence;
if the uplift is carried entirely by the target lag, the honest finding is
``neighbor activity predicts,'' reported as provisional. no-wlag and
wlag-only bound the contributions of point features and hydrology
respectively.

\paragraph{Target-lag leakage and its control.}
\textit{wlag\_nfip\_claims} is the spatial lag of the prediction target.
Computed once over the full dataset, it carries test-fold target values
into training: under spatially-blocked cross-validation a unit's neighbors
may sit in the held-out fold, so a training unit's lag is built partly from
held-out labels. This inflates precisely the relational uplift the ablation
is meant to measure---the test of whether spatial structure is real would
itself be contaminated by target-lag leakage. We therefore recompute
\textit{wlag\_nfip\_claims} \emph{per fold} from training units only: for
each fold, a unit's lag is the row-standardized mean of its neighbors'
target values restricted to the training partition, left undefined when all
of a unit's neighbors fall in the test fold. The variants that contain the
column (full, wlag-only) require an adjacency graph and are refused at
run time if it is absent, rather than silently falling back to the leaked
pre-computed values; the variants that do not (no-wlag, and the
headline no-target-lag) are unaffected and run without adjacency. The
headline relational verdict is thus clean of target-lag leakage by
construction.

\paragraph{Residual caveats.} The non-target spatial lags
(\textit{wlag\_rainfall\_mm} and the static neighbor lags) remain
full-dataset quantities; as lags of \emph{inputs} rather than the target,
they carry the ordinary spatial-smoothing dependence common to all spatial
features, not answer leakage, and are not recomputed per fold. A
residual-based diagnostic (\textit{spatial\_lag\_residual\_R0}) requiring
out-of-fold R0 predictions is specified but deferred to the R0 results
lock.

% ============================================================
% APPENDIX B — Out-of-distribution clustering probe (Oahu, s036)
% Real measured numbers (NOT \PH). Framed as instrument-transfer, not a
% sixth benchmark scenario; honest inconclusive-negative, no "unbiased" claim.
% ============================================================
\section{Out-of-Distribution Probe: Clustering on a Non-CONUS Substrate (Oahu)}
\label{app:oahu}

To test whether the clustering instrument (\textit{kappa\_spatial},
\S\ref{app:protocol-substrate}) transfers off the CONUS substrate, we apply
it to a structurally different setting: Oahu, Hawaii. This is a separate
experiment from the main benchmark---a different model (the Floodcaster
HAZUS damage engine, not the s035 representation ladder), a different region
(island topology, no continental adjacency), and a single decision geometry
(clustering only). It is an instrument-transfer probe, \emph{not} a sixth
benchmark scenario, and it carries no claim about the other five geometries.

\begin{figure*}[t]
  \centering
  \includegraphics[width=\textwidth]{figures/oahu_h4_results.pdf}
  \caption{Out-of-distribution clustering probe on Oahu (experiment s036,
    Floodcaster/HAZUS engine). The H4 clustering hypothesis \emph{fails}:
    residual autocorrelation is negative (Moran's $I=-0.358$, dispersion
    not clustering; \textit{kappa\_spatial}$=0.352$). At $n{=}17$ ZCTAs,
    3.5\% inundation, and a cumulative (non-event-matched) NFIP reference,
    this is an \emph{inconclusive negative}---not evidence that the damage
    model is geographically unbiased. The finding is a scope condition on
    the instrument (see text).}
  \label{fig:oahu}
\end{figure*}

% Companion figure (Floodcaster HAZUS depth--damage engine).
% Generated by: render_hazus_depth_damage.py --upload
\begin{figure*}[t]
  \centering
  \includegraphics[width=\textwidth]{figures/hazus_depth_damage_curves.pdf}
  \caption{FEMA HAZUS depth--damage functions used by the Floodcaster
    engine for the Oahu probe (\emph{not} the benchmark's HistGBDT/Ridge
    model). Panel~C shows Oahu observed damage by depth band.}
  \label{fig:hazus}
\end{figure*}

\paragraph{Setup.} A 10-year coastal-surge inundation raster with reef
attenuation (\texttt{Oahu\_10\_withReef}) drives HAZUS damage estimates for
47{,}983 buildings; 1{,}676 (3.5\%) are inundated (median depth 1.74\,ft,
94\% single-family slab-on-grade). Predicted losses are aggregated to the 17
of 20 Hawaii ZCTAs that contain buildings, compared against \emph{cumulative}
NFIP historical claims, normalized to $[0,1]$, and the absolute residual is
taken. Global Moran's $I$ is computed on the 17-ZCTA queen-contiguity graph.

\begin{table}[h]
\centering
\caption{Oahu clustering probe vs.\ the CONUS substrate. The probe returns
  negative (dispersed), not positive (clustered), residual autocorrelation.}
\label{app:tab:oahu}
\small
\begin{tabular}{@{}ll@{}}
\toprule
\textbf{Quantity} & \textbf{Value} \\
\midrule
Buildings / inundated        & 47{,}983 / 1{,}676 (3.5\%) \\
ZCTAs modeled                & 17 (of 20 HI) \\
Moran's $I$ (residuals)      & $-0.358$ \\
Null expectation $E[I]$      & $-0.0625$ \\
\textit{kappa\_spatial}      & $0.352$ (below the 0.5 neutral point) \\
Reference                    & cumulative NFIP (not event-matched) \\
\bottomrule
\end{tabular}
\end{table}

\paragraph{Result and interpretation.} Residual autocorrelation is
\emph{negative} ($I=-0.358$ vs.\ null $-0.0625$): adjacent ZCTAs tend to have
\emph{dissimilar} error magnitudes, the opposite of clustering. No positive
geographic clustering of residuals is detected, so the hypothesis that
\textit{kappa\_spatial} would flag reef/coastal discontinuities is not
supported on this substrate. We do \emph{not} read this as evidence that the
damage model is geographically unbiased: the same dispersed value is the
expected signature of an \emph{underpowered or confounded} test. Four
conditions make the negative result inconclusive rather than reassuring:
(i) at 3.5\% inundation, most ZCTAs carry near-zero predicted loss, so the
few wet ZCTAs alternate with dry neighbors and induce anti-correlation by
construction; (ii) ZCTA aggregation cross-cuts the reef line rather than
aligning with it, so a real coastal discontinuity is averaged away;
(iii) the reference is cumulative NFIP history, not the surge event, so the
residual measures model-vs-history mismatch, not model error; and
(iv) $n{=}17$ ZCTAs gives Moran's $I$ low power, and the sign could be noise.

\paragraph{What the probe contributes.} The value is a \emph{scope condition
on the instrument}, not a verdict on the model: \textit{kappa\_spatial} is
diagnostic for the clustering geometry only when wet-cell density and an
event-matched reference are adequate; under extreme inundation sparsity and a
historical reference it returns dispersion that must not be interpreted as
absence of geographic bias. For the Oahu substrate the clustering geometry is
therefore reported as \textsc{evidence-insufficient}---a clean demonstration
that the diagnostic, like any instrument, has an operating envelope the
benchmark should declare.
